% \numberwithin{equation}{section}

\chapter{Conclusions}

En aquest Treball de Fi de Grau s'han analitzat els modes normals d'oscil·lació d'un model de protuberància solar format per dos medis homogenis que representen el plasma de la protuberància i de la corona.
S'ha suposat que el camp magnètic és uniforme i transversal a la protuberància i que la gravetat es pot menysprear.
S'ha observat que les pertorbacions estacionàries d'aquest estat d'equilibri depenen tant de la constant d'acoblament, $k_z$,
com dels paràmetres del sistema, en concret, de la velocitat del so, $c_s$, i la velocitat d'Alfvén, $v_A$,
tant en l'interior de la protuberància com en l'exterior, en la corona.

La subdivisió del problema en les Eqs.~(\ref{eq.3.30}), (\ref{eq.3.31}) i (\ref{eq.3.32}) condueix a dos tipus de solucions,
els modes d'Alfvén, molts semblants als d'una corda amb major densitat en el centre, i els modes ràpids i lents.
En el cas d'acoblament nul, $k_z = 0$, aquests darrers presenten la mateixa estructura que els modes d'Alfvén i, en conseqüència, que els de la corda.

Els autovalors han estat calculats de manera correcta per Dedalus.
Aquest fet s'ha comprovat mitjançant el càlcul de l'error relatiu, $\eta$, respecte a les solucions analítiques amb forma de funcions trigonomètriques,
d'acord amb l'article de \textcite{oscillation_modes_II}.
L'error relatiu és molt menor a la unitat en els primers 100 autovalors, tot i que aquest creix significativament quan el nombre de l'autovalor augmenta. 
La situació observada també té lloc en el càlcul de les autosolucions d'una corda, i es deu a la manera de treballar de Dedalus,
que proporciona una major precisió quan les autofuncions tenen pocs extrems en la direcció $\hat{x}$, és a dir, quan la freqüència és baixa.
Amb un nombre, N, major de funcions base pel càlcul d'autovalors s'observa la mateixa tendència, conservant les proporcions.

Les condicions del sistema juguen un paper clau a l'hora de determinar la constant d'acoblament, $k_z$.
En conseqüència, els modes normals resultants es veuen directament afectats per els condicions del sistema, llavors els modes ràpids i lents també.
Aquest fet es recflecteix clarament en els encreuaments evitats.

L'estudi de les autofuncions és imprescindible a l'hora d'entendre el que succeeix en els encreuaments evitats dels modes normals,
és a dir, en els acoblaments dels modes ràpids i lents.
En el cas dels modes ràpids, el mode no dominant, $\tilde{v}_{x}$, és aquell que presenta una direcció paral·lela al camp magnètic del sistema, $\bm{B}$,
mentre que el mode majoritari, $v_z$, és perpendicular al camp i paral·lel al nombre d'ona o constant d'acoblament, $k_z$.
Aquests sempre presenten paritats oposades, com s'ha demostrat en les diferents figures del treball.
Els modes lents en canvi, demostren un comportament oposat. Els seus modes no dominants, $v_z$, són paral·lels a l'acoblament,
mentre que els predominants, $\tilde{v}_{x}$, són aquells perpendiculars al vector d'ona i paral·lels al camp.

S'ha demostrat el canvi de dominància d'un mode a un altre en el sí d'un encreuament.
A més, s'ha estudiat l'evolució dels modes fonamentals híbrids per un ampli rang de valors de la constant d'acoblament, $k_z$.
D'aquesta manera es pot apreciar com els modes lents no es veuen molt afectats per un acoblament major,
mentre que els ràpids en pateixen les conseqüències de manera més notable, tant en les autofuncions dominants com en les no dominants.

Aquest és un estudi introductori d'acord amb les limitacions de temps i dificultat que suposa un treball de fi de grau.
Una manera d'ampliar aquest projecte i augmentar la seva complexitat consisteix en realitzar manco aproximacions al sistema de la Fig.~\ref{fig:sistema},
com per exemple incloure una component en la direcció $\hat{y}$ al camp magnètic en l'equilibri, fet per \textcite{oscillation_modes_III}.
D'aquesta manera les Eqs.~(\ref{eq.3.21})~i~(\ref{eq.3.23}) estan compostes per termes extres tals que lliguen dites equacions.
L'anàlisi dels modes ràpids, lents i d'Alfvén, tots acoblats entre sí excepte per $k_z = 0$, que sorgeixen d'aquest cas més general es deixa per un futur treball.